% Principal: GDZ PPN236018647, native 18–19, printed 11–12.
\SourceLeaf{18}\NoSourceHead
\pdfbookmark[1]{Einige asymptotische Gesetze der Zahlentheorie}{asymptotische}
\begin{center}\large NACHLASS.\end{center}
\HistoricTitle{EINIGE ASYMPTOTISCHE GESETZE DER ZAHLENTHEORIE.}
\SourceGroup{I}
\HistoricCitation{[Handschriftliche Eintragung in dem Buche:] \textsc{Johann Carl Schulze}, Neue und erweiterte Sammlung logarithmischer . . . . Tafeln. I, Berlin 1778; [von \textsc{Gauss}' Hand] \textbf{Gauß. 1791.}}

\noindent\small [Auf der Rückseite des letzten Blattes.]\normalsize
\NumberedItem{1}
\par\indent Primzahlen unter \(a\ (=\infty)\)
\[
\frac{a}{la}.
\]
\ClosingRule
\NumberedItem{2}
\par\indent Zahlen aus zwei Factoren
\[
\frac{lla.a}{la},
\]
\noindent (wahrsch.) aus 3 Factoren
\[
\frac{\frac12(lla)^2a}{la},\ .\,.\,.
\]
\noindent et sic in inf.
\ClosingRule
\vfill\begin{flushright}\small 2\textsuperscript{*}\end{flushright}
\clearpage
\SourceLeaf{19}\ArithmeticHead{12}{L}{ARITHMETIK.\quad NACHLASS.}
\NumberedItem{3}
\[
1+\frac12+\frac13+\frac15+\frac17+\frac1{11}+.\,.\,.+\frac1x
=\text{(pro }x\text{ inf.) }llx+V.
\]
\noindent \(V\) esse Const. suspicor ac prope \(1,266\ .\,.\,.\)
\begin{flushright}1796 Mai.\end{flushright}
\ClosingRule
\NumberedItem{4}
\par\indent Zahlen die keine gleichen F[actoren] haben
\[
\frac{a}{1+\frac14+\frac19+.\,.\,.},
\]
\noindent die höchstens 2 gl. F. h.
\[
\frac{a}{1+\frac18+\frac1{27}+.\,.}
\]
\noindent die höchstens 3 gl. F. h.
\[
\frac{a}{1+\frac1{16}+\frac1{81}+.\,.\,.},
\]
\noindent et sic in inf.
\ClosingRule
\NumberedItem{5}
\[
\frac21.\frac32.\frac54.\frac76.\,.\,.\frac{x}{x-1}
=(x\text{ inf})\,a.lx
\]
\noindent \(a\) Constans ac prope \(1,874\).
\ClosingRule
\bigskip\noindent\small [Auf der Rückseite des hinteren Schutzblattes]\normalsize
\NumberedItem{6}
\par\indent Numerus factorum usque ad \(n\)
\[
(lln+1)n.
\]
\clearpage
